\section{Data}\label{sec:data}

We demonstrate the \texttt{microimpute} methodology through an application to wealth imputation, transferring net worth from the Survey of Consumer Finances (SCF) onto the Current Population Survey (CPS).

\subsection{Survey of Consumer Finances}

The Survey of Consumer Finances (SCF), sponsored by the Federal Reserve Board, is a triennial survey providing detailed information on U.S. households' assets, liabilities, income, and demographic characteristics. Its dual-frame sample design includes a standard national area-probability sample and a list sample deliberately oversampling wealthy households to better capture the skewed wealth distribution \citep{barcelo2006imputation}. The SCF is a benchmark for wealth imputation research due to its detailed financial data and the known complexities arising from its design and the nature of wealth. Item nonresponse in public-use SCF datasets is addressed by the Federal Reserve through a multiple imputation approach that generates five complete datasets with different imputed values, using sequential regression-based procedures that incorporate range constraints, logical data structures, and empirical residuals to preserve the complex multivariate relationships inherent in wealth data \citep{kennickell1998multiple}. 

Specifically, we use the 2022 summarized SCF as our donor dataset. 

\subsection{Current Population Survey}

The Current Population Survey (CPS), conducted jointly by the U.S. Census Bureau and the Bureau of Labor Statistics, is a monthly survey of approximately 60,000 U.S. households that serves as the primary source of labor force statistics for the United States. The CPS uses a multistage probability-based sample designed to represent the civilian non-institutional population; on average, each sampled household represents approximately 2,500 households in the population. The Annual Social and Economic Supplement (ASEC) extends the core survey with detailed annual income data, including earnings, unemployment compensation, Social Security, pension income, interest, dividends, and other income sources. However, despite its comprehensive coverage of income and employment, the CPS does not collect information on household wealth, assets, or liabilities, which are key variables needed for wealth-based policy analysis. This omission motivates the need for statistical matching to transfer wealth information from the SCF onto the CPS.

\subsection{Comparative analysis and characteristics for imputation}

Imputing between SCF and CPS exemplifies the challenges of statistical matching, as it requires combining a specialized survey with detailed wealth data but limited sample size with a large representative survey that lacks wealth measures entirely. The SCF-to-CPS imputation is directly relevant to U.S. policy microsimulation, where researchers require comprehensive microdata combining income, demographics, and wealth to analyze the distributional effects of tax and benefit reforms. By imputing wealth onto the CPS, analysts can examine wealth distributions in a sample roughly thirteen times larger than the SCF alone, enabling finer demographic disaggregations and more precise subgroup analyses. Moreover, the heavy-tailed nature of wealth distributions and the different sampling designs of these surveys create methodological challenges that make this an informative test case for comparing imputation methods.